% Chapter 16: The Willmore Conjecture
\let\LocalMacro\relax

\chapter{The Willmore Conjecture}

\section{The Problem}
For a closed immersed surface $\Sigma \subset \mathbb{R}^3$, the \emph{Willmore energy} is
\[
\mathcal{W}(\Sigma) = \int_\Sigma H^2 \, dA
\]
where $H = \frac{1}{2}(k_1 + k_2)$ is the mean curvature.

\begin{conjecture}[Willmore, 1965]
For any immersed torus $T^2 \subset \mathbb{R}^3$,
\[
\mathcal{W}(T^2) \ge 2\pi^2.
\]
Equality holds for the Clifford torus in $S^3$ (stereographically projected to $\mathbb{R}^3$).
\end{conjecture}

\begin{historicalbox}
Willmore conjectured this in 1965. It became a central problem in geometric analysis, connecting conformal geometry, integrable systems, and algebraic geometry. The conjecture resisted all attempts for 47 years.
\end{historicalbox}

\section{Mathematical Theory}


\subsection{Prerequisites}
This chapter assumes familiarity with:
\begin{itemize}
\item Differential geometry of surfaces (first/second fundamental forms, mean curvature)
\item Variational calculus and geometric measure theory
\item Basic complex analysis (Riemann surfaces, uniformization)
\item Conformal geometry
\end{itemize}

\subsection{The Willmore Functional and Conformal Invariance}
The Willmore energy is conformally invariant in $\mathbb{R}^3$:
\[
\mathcal{W}(\Sigma) = \int_\Sigma (H^2 - K) \, dA + \int_\Sigma K \, dA = \int_\Sigma (H^2 - K) \, dA + 2\pi\chi(\Sigma)
\]
For a torus ($\chi=0$), $\mathcal{W}(\Sigma) = \int (H^2 - K) \, dA$. The integrand $H^2 - K = \frac{1}{4}(k_1 - k_2)^2 \ge 0$ is conformally invariant.

\subsection{The Clifford Torus}
The Clifford torus in $S^3$ is
\[
T_{\text{Cl}} = \left\{ (z_1, z_2) \in \mathbb{C}^2 : |z_1| = |z_2| = \frac{1}{\sqrt{2}} \right\}.
\]
It is minimal in $S^3$, has constant mean curvature in $\mathbb{R}^3$ (after stereographic projection), and $\mathcal{W}(T_{\text{Cl}}) = 2\pi^2$.

\section{The Discovery}

\subsection{Marques and Neves \cite{marquesneves2014}}

\begin{theorem}[Marques-Neves \cite{marquesneves2014}]
The Willmore conjecture is true: $\mathcal{W}(T^2) \ge 2\pi^2$ for any immersed torus $T^2 \subset \mathbb{R}^3$, with equality only for the Clifford torus (up to conformal transformations).
\end{theorem}

Their proof used:
\begin{enumerate}
\item \textbf{Almgren-Pitts min-max theory}: Construct minimal surfaces in $S^3$ via sweepouts.
\item \textbf{The canonical family}: For a given torus $\Sigma$, construct a family of surfaces in $S^3$ by conformal transformations, whose maximum area gives a lower bound on $\mathcal{W}(\Sigma)$.
\item \textbf{The index bound}: Show that any minimal surface in $S^3$ with area $< 2\pi^2$ must be the Clifford torus.
\item \textbf{The free boundary case}: Extend to surfaces with boundary using free boundary min-max.
\end{enumerate}

\begin{highlightbox}
The proof is a tour de force of geometric measure theory: it combines Almgren-Pitts min-max, conformal geometry, and the classification of minimal surfaces in $S^3$ to solve a 47-year-old problem.
\end{highlightbox}

\section{Impact}
\begin{itemize}
\item \textbf{Geometric analysis}: Almgren-Pitts min-max became the standard tool for Willmore-type problems.
\item \textbf{Integrable systems}: Willmore surfaces are related to soliton equations; the Clifford torus corresponds to a finite-gap solution.
\item \textbf{Conformal geometry}: The Willmore energy is the natural energy for conformal immersions.
\item \textbf{Physics}: Willmore energy appears in membrane elasticity (Helfrich energy).
\end{itemize}

\section{Exercises}

\begin{exercise}[Basic]
Compute the Willmore energy of the Clifford torus in $\mathbb{R}^3$ (stereographically projected from $S^3$).
\end{exercise}

\begin{exercise}[Intermediate]
Show that the Willmore energy is conformally invariant in $\mathbb{R}^3$. Why does this fail in higher dimensions?
\end{exercise}

\begin{exercise}[Advanced]
Explain how Almgren-Pitts min-max produces a minimal surface in $S^3$ from a given sweepout of a torus.
\end{exercise}

\section{Timeline}

\begin{tabularx}{\linewidth}{lX}
\toprule
\textbf{Year} & \textbf{Event} \\ \midrule
1965 & Willmore conjectures $2\pi^2$ lower bound for tori \\ 
1993 & Li-Yau proves $\mathcal{W} \ge 2\pi^2$ for non-embedded tori \\ 
2005 & Ros proves Willmore conjecture for tori of revolution \\ 
2012 & Marques-Neves prove Willmore conjecture \\
2014 & Marques-Neves win Ramanujan Prize \\ \bottomrule
\end{tabularx}

\section{Citations}

\begin{enumerate}
\item Marques, F. C., \& Neves, A. (2014). ``Min-max theory and the Willmore conjecture.'' Annals of Mathematics, 179(2), 683--782.
\item Willmore, T. J. (1965). ``Note on embedded surfaces.'' Anais da Academia Brasileira de Ciências, 37, 475--478.
\item Simon, L. (1993). ``Existence of surfaces minimizing the Willmore functional.'' Communications in Analysis and Geometry, 1(2), 281--326.
\end{enumerate}
